\question （西安交通大学，2004年）某机浮点数格式为：数符1位、阶符1位、阶码5位、尾数9位（共16位）。若机内采用阶移尾补规格化浮点数表示，那么它能表示的最小负数为（
）
\par\fourch{\textcolor{red}{}}{}{}{}
\begin{solution}当阶码为5位时，移码所能表示的最大真值与补码是一样的，即31，故可以排除B和D。上面讲过，当尾数采用补码时，最大的规格化数是-1（1.000000000），故它能表示的最小负数为\includegraphics[width=0.34375in,height=0.15625in]{texmath/ea62b55Cdpi7B3507D-25E7B317D}。
\end{solution}
